\documentstyle[% 


righttag% 


,11pt% 


,amssymb% 


,russian%               remove in Eng. 


,izvru%                 Set izven% in Eng. 


% ,izvru-ks             for brief communications 


]{amsart} 


 


%       Math definitions 


 


\newcommand{\tts}[1]{} 


 


\theoremstyle{plain} 


\theorembodyfont{\it} 


\newtheorem{thmnonum}{\hskip\parindent Теорема} 


\newtheorem{lemnonum}{\hskip\parindent Лемма} 


\renewcommand{\thethmnonum}{} 


\renewcommand{\thelemnonum}{} 


\theoremstyle{definition} 


\theorembodyfont{\rm} 


\newtheorem{defnnonum}{\hskip\parindent Определение} 


\renewcommand{\thedefnnonum}{} 


 


 


%\hoffset=-15mm%                  For Eng. 


  


\setcounter{page}{7} 


\god{1998} 


\nomer{10~(437)} %                        Remove in Eng. 


%\nomer{Vol.\,42, No.\,10}%               Set in Eng. 


%\str{\pageref{begin}--\pageref{end}}%  Set in Eng. 


\udk{517.956} 


\author{\it М.К.\,Бугир} 


% NB:   ^^ remove in Eng. 


\title{Оценки условной устойчивости\\ 


двухконтурной краевой задачи в круге} 


 


\date{} 


% \pagestyle{myheadings}%         Set in Eng. 


 


\pagestyle{plain} %              Remove in Eng. 


\begin{document} 


 


% \thispagestyle{plain}%          Set in Eng. 


 


\maketitle 


\label{begin} 


 


{\bf 1.} В \cite{1}--\cite{3} рассматривалась разноконтурная краевая 


задача для полигармонического уравнения 


\begin{equation} 


\Delta^n u+c_1\Delta^{n-1} u+\dotsb+c_n u=0 


\tag{1${}_{2n}$} 


\end{equation} 


с оператором Лапласа $\Delta$ и постоянными коэффициентами 


$c_1,\dotsc,c_n$, когда краевые условия задавались на разных контурах 


\begin{equation} 


u(x)\big|_{\Gamma_1}=\varphi_1(x),\dotsc,\ \; 


u(x)\big|_{\Gamma_n}=\varphi_n(x), \tag{2${}_{2n}$} 


\end{equation} 


где замкнутые контуры $\Gamma_k$ ограничивают области $D_k$, причем 


$D_1\subset D_2\subset\dotsb\subset D_n\subseteq D$. Кроме того, 


как показано в \cite{2}, \cite{3}, результаты практически не меняются, 


если вместо оператора Лапласа $\Delta$ взять оператор Лапласа--Бельтрами 


$L$ и рассмотреть краевую задачу (2) для соответствующих уравнений в 


пространствах постоянной кривизны. Отметим, кроме того, что такие задачи 


являются одним из возможных обобщений многоточечных задач для 


обыкновенных дифференциальных уравнений и возникают в самых 


разнообразных прикладных проблемах, связанных, например, с 


геологическими, геофизическими и другими наблюдениями \cite{4}, когда 


измерение экспериментальных данных осуществляется не только на границе 


области, но и внутри нее. 


 


Ниже будет показано, что задача ($1_{2n}$), ($2_{2n}$) некорректная, 


поэтому 


важную роль играют оценки условной устойчивости задачи, классическим 


результатом в этом смысле можно считать оценки для задачи аналитического 


продолжения Адамара (\cite{4}, с.\,26--27). В данной работе мы также 


остановимся на оценках условной устойчивости для задачи ($1_{2n}$), 


($2_{2n}$) при $n=2$, т.\,е.~рассмотрим уравнение четвертого порядка 


\begin{equation} 


K_4u=\Delta^2u+c_1\Delta u+c_2u=0,\tag{1${}_4$} 


\end{equation} 


когда контуры $\Gamma_1$ и $\Gamma_2$ имеют вид концентрических 


окружностей радиусов $R_1$ 


и $R_2$, $R_2>R_1$. (В дальнейших обозначениях (1), (2) нижний индекс 


4 будем опускать, если это не вызывает недоразумений.) 


 


Уравнение (1) через корни квадратного уравнения 


\addtocounter{equation}{2}


\begin{equation} 


\lambda^2+c_1\lambda+c_2=0 


\end{equation} 


можно переписать в виде суперпозиции двух операторов второго порядка 


$K_4u=(\Delta-\lambda_2)(\Delta-\lambda_1)u=0$ 


или эквивалентной системы уравнений второго порядка 


\begin{equation} 


\begin{cases} 


(\Delta-\lambda_1)u=v, \\ 


(\Delta-\lambda_2)v=0. 


\end{cases} 


\end{equation} 


В силу (4) произвольное решение (1) представляется в виде 


\begin{equation} 


u=u_1+\wt v, 


\end{equation} 


где $u_1$ --- произвольное решение уравнения 


\begin{equation} 


K_{i,2}u_i=(\Delta-\lambda_i)u_i=0\quad (i=1,2) \tag{6${}_i$} 


\end{equation} 


при $i=1$, а $\wt v$ --- частное решение неоднородного уравнения 


\begin{equation} 


(\Delta-\lambda_1)u=v.\tag{4${}_1$} 


\end{equation} 


 


В \cite{5} представление произвольного решения (5) в зависимости от вида 


корней уравнения (3) выписано более точно. В частности, если эти корни 


разные, то 


\begin{equation} 


u=u_1+u_2,\tag{5$'$} 


\end{equation} 


где $u_i$ --- произвольные решения уравнений (6${}_i$) ($i=1,2$) и 


\begin{equation} 


u=u_0+r\frac{\partial u_1}{\partial r}\tag{5$''$} 


\end{equation} 


для кратного корня уравнения (3), где $u_0$ и $u_1$ --- произвольные 


решения уравнения $\Delta u-\lambda u=0$. 


 


\begin{defnnonum} 


Нетривиальное решение уравнения (1) будем называть колеблющимся в 


области $D$, если найдутся хотя бы два замкнутых контура $\Gamma_1$ и 


$\Gamma_2$, ограничивающие области $D_1$ и $D_2$, 


$D_1\subset D_2\subseteq D$, такие, что 


$$ 


u(x)\big|_{\Gamma_1}=0=u(x)\big|_{\Gamma_2},


$$ 


и неколеблющимся в противном случае. Уравнение (1) будем называть 


неколеблющимся, если все его решения неколеблющиеся.


\end{defnnonum} 


 


Данное определение колеблемости связано с порядком уравнения и легко 


обобщается на уравнения порядка $2n$. 


 


Будем предполагать, что корни уравнения (3) неотрицательные, что 


обеспечивает неколеблемость уравнения (1). 


 


Решение уравнения (1) для разных корней уравнения (3) в полярных 


координатах $(r,\theta)$ с помощью метода разделения переменных запишем 


через модифицированные функции Бесселя $I_m(t)$ следующим образом: 


\addtocounter{equation}{1}


\begin{equation} 


u(r,\theta)=\sum^\infty_{m=0}\sum^2_{j=1}c_{m,j}I_m 


(\sqrt{\lambda_j} r\,)e^{im\theta},\quad i=\sqrt{-1}. 


\end{equation} 


 


Для вычисления коэффициентов Фурье представления (7) получим систему 


алгебраических уравнений относительно неизвестных коэффициентов 


$c_{m,j}$ ($j=1,2$; $m=0,1,2,\dotsc$) 


$$ 


\sum^2_{j=1}c_{m,j}I_m(\sqrt{\lambda_j}\,r)=\varphi_{m,l}\quad 


(l=1,2), 


$$ 


где правые части суть коэффициенты разложения Фурье контурных условий 


(2). Единственность решения последней алгебраической системы уравнений 


зависит от значения определителя 


$$ 


\Delta_m= 


\begin{vmatrix} 


I_m(\sqrt{\lambda_1}R_1) & I_m(\sqrt{\lambda_2}R_1) \\ 


I_m(\sqrt{\lambda_1}R_2) & I_m(\sqrt{\lambda_2} R_2) 


\end{vmatrix} 


\quad (m=0,1,2,\dotsc). 


$$ 


 


Очевидно, необходимым и достаточным условием единственности решения 


двухконтурной краевой задачи (1), (2) будет 


\begin{equation} 


\Delta_m\ne 0\quad (m=0,1,2,\dotsc). 


\end{equation} 


Последнее условие обеспечивается, как уже упоминалось, условиями 


\begin{equation} 


\lambda_1\ge 0,\quad \lambda_2\ge 0. 


\end{equation} 


 


Используем обозначения работ \cite{1}--\cite{3} $R_2/R_1=\tau$, 


$\delta_m(t)=I_m(t)/I_m(\tau t)$, 


$\sigma_i=\sqrt{\lambda_i}\,R_1$ ($i=1,2$). Если в 


ряде (7) сгруппируем выражения, стоящие при коэффициентах Фурье 


$\varphi_{m,1}$ и $\varphi_{m,2}$, то он перепишется в виде 


\begin{equation} 


u(r,\theta)=\sum^\infty_{m=0}\sum^2_{j=1}B_{m,j}(r) 


e^{im\theta}\varphi_{m,j}, 


\end{equation} 


где 


$$ 


B_{m,1}(r)= 


\frac{ 


\begin{vmatrix} 


\displaystyle 


\frac{I_m(\sqrt{\lambda_1}r)}{I_m(\tau\sigma_1)} & 


\displaystyle 


\frac{I_m(\sqrt{\lambda_2}r)}{I_m(\tau\sigma_2)} \\ 


1 & 1 


\end{vmatrix} 


} 


{\delta_m(\sigma_1)-\delta_m(\sigma_2)} 


,


\qquad 


% 


B_{m,2}(r)= 


\frac{ 


\begin{vmatrix} 


\delta_m(\sigma_1) & \delta_m(\sigma_2) \\ 


\displaystyle 


\frac{I_m(\sqrt{\lambda_1}r)}{I_m(\tau\sigma_1)} & 


\displaystyle 


\frac{I_m(\sqrt{\lambda_2}r)}{I_m(\tau\sigma_2} 


\end{vmatrix} 


} 


{\delta_m(\sigma_1)-\delta_m(\sigma_2)}. 


$$ 


\medskip 


 


{\bf 2.} Рассмотрим вопрос об условной устойчивости задачи (1), (2) при 


дополнительных условиях на решения уравнения (1). Как указано в работе 


\cite{6}, одним из таких естественных условий является 


\begin{equation} 


\frac{2}{\pi}\int^{2\pi}_0\bigg[ 


\biggl(\frac{\partial u}{\partial r}\biggr)^2+ 


\biggl(\frac{\partial u}{\partial\theta}\biggr)^2\bigg]\bigg|_{r=R_2} 


d\theta<E, 


\end{equation} 


выражающее ограниченность энергии. 


 


Некорректность двухконтурной задачи (1), (2) нетрудно заметить, если 


задать значения на контурах $\varphi_2(x)=0$, 


$\varphi_1(x)=\varepsilon{}e^{im\theta}$. 


 


Рассмотрим пространства 


\begin{gather} 


L_2[0,2\pi]\quad 


\text{с нормой}\quad 


\|\varphi(\theta)\|=\frac{2}{\pi}\int^{2\pi}_0 


\varphi^2(\theta)d\theta=\sum^\infty_{k=0}\varphi^2_k, \\ 


H_2\quad 


\text{с нормой}\quad 


\|u(r,\theta)\|^2_{H_2}=\max_{0\le r\le R_2}\frac{2}{\pi} 


\int^{2\pi}_0 u^2(r,\theta)d\theta=\max_{0\le r\le R_2} 


\sum^\infty_{k=0}u^2_k(r),\tag{12$'$} 


\end{gather} 


где $\varphi_k$ и $u_k(r)$ --- коэффициенты разложения Фурье 


соответствующих функций. 


Рассмотрим решение 


$u_{m,\varepsilon}(x)=\varepsilon{}B_{m,1}(r)e^{im\theta}$, 


удовлетворяющее указанным выше условиям. В 


пространстве $L_2[0,2\pi]$ для произвольного $m$ и произвольно малого 


$\varepsilon>0$\; $\|\varphi_1\|_{L_2[0,2\pi]}=\varepsilon$, 


т.\,е.~произвольно мала. С другой стороны, в силу оценки для функции 


$\delta_m(t)$ \cite{2} 


$$ 


|B_{m,1}(r)|= 


\frac{ 


\displaystyle\bigg| 


\frac{I_m(\sqrt{\lambda_1}r)}{I_m(\sqrt{\lambda_1}R_2)}- 


\frac{I_m(\sqrt{\lambda_2}r)}{I_m(\sqrt{\lambda_2}R_2)} 


\bigg|} 


{|\delta_m(\sqrt{\lambda_1}R_1)-\delta_m(\sqrt{\lambda_2}R_2)|} 


\ge[\max\big(\delta_m(\sqrt{\lambda_1}R_1)-\delta_m 


(\sqrt{\lambda_2}R_2)\big)]^{-1}=\wt A\tau^m. 


$$ 


Последнее выражение за счет $m$ можно выбрать сколь угодно большим. 


Следовательно, некорректность двухконтурной задачи проявляется в том, 


что нет непрерывной зависимости по норме (12$'$) от исходных данных на 


внутреннем контуре $\Gamma_1$. 


 


\begin{note} 


Нетрудно проверить, что априорные оценки условной устойчивости 


двухконтурной задачи (1), (2) тесным образом связаны с условиями 


единственности аналитического продолжения функции $\wt v(x)$ с области 


$D_1$ на всю область $D_2$. Для этого достаточно записать систему 


интегральных уравнений, равносильную системе уравнений (4),


\begin{equation} 


\begin{split} 


u(P_0)&=\idotsint\limits_{D_2}(\lambda_1u+w)G_2(P,P_0)dP+ 


\idotsint\limits_{\Gamma_2}u(P) 


\frac{\partial G_2(P,P_0)}{\partial n_P}dP, \\ 


% 


w(P_0)&=\idotsint\limits_{D_1}(\lambda_2w G_1(P,P_0)dP+ 


\idotsint\limits_{\Gamma_1}w(P) 


\frac{\partial G_1(P,P_0)}{\partial n_P}dP,


\end{split} 


\end{equation} 


где через $G_i(P,P_0)$ обозначены функции Грина задачи с условиями 


Дирихле для 


уравнения Лапласа в областях $D_i$ ($i=1,2$). Как известно \cite{7}, 


второе уравнение (13) всегда имеет единственное решение в области $D_1$, 


и если 


его единственным образом продолжить на всю область, то первое уравнение 


также будет иметь единственное решение. В связи с этим сформулируем 


дополнительные условия типа условий (11), обеспечивающие условную 


устойчивость задачи (1), (2), через частное решение $\wt v(x)$ 


неоднородного уравнения (4${}_1$). 


\end{note} 


 


\begin{thm} 


Пусть частное решение $\wt v(r,\theta)$ неоднородного уравнения $(4_1)$ 


удовлетворяет априорной оценке 


\begin{equation} 


\frac{2}{\pi}\int^{2\pi}_0\wt v^{\,2}(R_2,\theta)d\theta<M^2 


\end{equation} 


и 


$$ 


\|\varphi_i(\theta,R_i)\|_{L_2[0,2\pi]}\le\varepsilon_i\quad 


(i=1,2). 


$$ 


Тогда задача $(1)$, $(2)$ условно устойчива и имеет место оценка 


\begin{equation} 


\frac{2}{\pi}\int^{2\pi}_0 u^2(r,\theta)d\theta\le 2\bigg\{ 


\varepsilon_1+ 


\frac{I_k(\sqrt{\lambda_2}r)}{I^2_k(\sqrt{\lambda_2}R_2)} 


M^2(\lambda_2-\lambda_1)^2\bigg\}, 


\end{equation} 


где число $k$ определяется из неравенства 


\begin{equation} 


0\le k+(1/2)\ln 2(\varepsilon_1+\varepsilon_2)M^{-2}/ 


(\ln\tau\beta/\ln\tau), 


\end{equation} 


если корни уравнения $(3)$ разные, a для кратных корней уравнения $(3)$ 


\begin{equation} 


\frac{2}{\pi}\int^{2\pi}_0u^2(r,\theta)d\theta\le 2\bigg\{\varepsilon_1+ 


M^2\bigg[ 


\frac{r I'_k(\sqrt{\lambda}r)}{R_2I'(\sqrt{\lambda}R_2)} 


\bigg]\bigg\}, 


\end{equation} 


где $k$ удовлетворяет уравнению 


\begin{equation} 


M\frac{R_1}{R_2} 


\frac{I'_k(\sqrt{\lambda}R_1)}{I'_k(\sqrt{\lambda}R_2)} 


=\sqrt{2(\varepsilon_1+\varepsilon_2)}\,. 


\end{equation} 


\end{thm} 


 


\begin{pf} 


В силу представления (5) решение задачи (1), (2) можно разбить на две 


задачи: задачу с условиями Дирихле 


$$ 


u_1\big|_{\Gamma_2}=\varphi_2(x)


$$ 


для уравнения (6${}_1$) и двухконтурную задачу для решения $\wt v(x)$ 


$$ 


\wt v\big|_{\Gamma_2}=0,\quad 


\wt v\big|_{\Gamma_1}=\varphi_1-u_1\big|_{\Gamma_1}. 


$$ 


 


Из-за неколеблемости уравнений (6${}_i$) имеет место принцип максимума 


для них, поэтому в круге $r\le R_2$ выполняются неравенства 


$$ 


\frac{2}{\pi}\int^{2\pi}_0u^2_i(r,\theta)d\theta<\varepsilon_2, 


\quad r\in[0,R_2]\ \; (i=1,2)


$$ 


при условии $\|\varphi_2\|_{L_2}<\varepsilon_2$. Так как 


$-\wt v^{\,2}\le 2(\varphi^2_1+u^2_1\big|_{\Gamma_1})$ на контуре 


$\Gamma_1$, то 


\begin{equation} 


\frac{2}{\pi}\int^{2\pi}_0\wt v^{\,2}(R_1,\theta)d\theta 


<2(\varepsilon_1+\varepsilon_2). 


\end{equation} 


Разделяя переменные в полярной системе координат, решения однородных 


уравнений (6${}_i$) можно переписать в виде 


$$ 


u_i(r,\theta)=\frac{a_0}{2}+\sum^\infty_{k=1} 


I_k(\sqrt{\lambda_i}r)a_k e^{ik\theta}. 


$$ 


Последнее выражение можно рассматривать как разложение в ряд Фурье на 


интервале $[0,2\pi]$, поэтому 


$$ 


\frac{2}{\pi}\int^{2\pi}_0u^2_i(r,\theta)d\theta= 


\frac{A^2_0}{4}+\sum^\infty_{k=1}A^2_k I^2_k 


(\sqrt{\lambda_i}r). 


$$ 


Перепишем частное решение неоднородного уравнения (4${}_1$) \cite{8} 


$$ 


\wt v(r,\theta)=\sum^\infty_{k=1}I_k 


(\sqrt{\lambda_2}r)e^{ik\theta} 


\frac{A_k}{\lambda_2-\lambda_1}+\frac{A_0}{2},


$$ 


когда $\lambda_1\ne\lambda_2$, и 


$$ 


\wt v(r,\theta)=r\sum^\infty_{k=1}I'_k 


(\sqrt{\lambda}r)e^{ik\theta}A_k+r\frac{A_0}{2}, 


$$ 


если $\lambda_1=\lambda_2=\lambda$. (В дальнейшем ограничимся только 


первым случаем, когда $\lambda_1\ne\lambda_2$, 


т.\,к.~во втором случае рассуждения аналогичные с заменой функции 


$I^2_k(\lambda_2r)$ на $r^2I'_k{}^2(\lambda_2r)$.) 


 


Учитывая ортогональность множества функций $\{e^{ik\theta}\}$ в 


пространстве $L_2[0,2\pi]$ и 


условия (14) и (19) при разных корнях уравнения (3), получим систему 


неравенств-ограничений 


\begin{equation} 


\begin{cases} 


\displaystyle 


\frac{1}{(\lambda_1-\lambda_2)^2}\sum^\infty_{k=0}A^2_k I^2_k 


(\sqrt{\lambda_2}R_1)\le2(\varepsilon_1+\varepsilon_2),\\ 


\displaystyle 


\sum^\infty_{k=0}A^2_k I^2_k(\sqrt{\lambda_2}R_2)\le 


M^2(\lambda_2-\lambda_1)^2. 


\end{cases} 


\end{equation} 


Неизвестные $A^2_k$ при каждом $r\in[0,R_2]$ 


следует подобрать так, чтобы функция 


$$ 


F(r)=\int_0^{2\pi} v^2(r,\theta)d\theta= 


\frac{A^2_0}{4}+\sum^\infty_{k=0}A^2_k I^2_k(\lambda_2r) 


$$ 


принимала наибольшее значение. Нетрудно заметить, что построенная выше 


задача является задачей линейного программирования, поэтому оптимальное 


значение достигается на некотором опорном решении. Следовательно, 


существуют такие $k_1<k_2$, что $A_{k_1}$, $A_{k_2}$ отличны от нуля, а 


остальные $A_j=0$, если 


оптимальное решение невырождено, и $A_k\ne 0$, если оно вырожденное. 


Используя 


обозначения $I^2_{k_j}(\sqrt{\lambda_2}R_i)=a_{ij}$, для определения 


неизвестных $A_{k_1}$ и $A_{k_2}$ получим систему уравнений 


\begin{equation} 


\begin{cases} 


A^2_{k_1}a_{11}+A^2_{k_2}a_{12}=2(\varepsilon_1+\varepsilon_2) 


(\lambda_2-\lambda_1)^2,\\ 


A^2_{k_1}a_{21}+A^2_{k_2}a_{22}=M^2(\lambda_2-\lambda_1)^2. 


\end{cases} 


\end{equation} 


Для вырожденного опорного решения утверждение теоремы почти очевидно, 


поскольку\linebreak $A_{k_1}=0$, 


$A^2_{k_2}=M^2(\lambda_2-\lambda_1)^2I^{-2}_{k_2} 


(\sqrt{\lambda_2}R_2)$ или наоборот, а число 


$k$ определяется из трансцендентного уравнения 


$$ 


M^2\frac{I^2_k(\sqrt{\lambda_2}R_1)} 


{I^2_k(\sqrt{\lambda_2}R_2)}=2(\varepsilon_1+\varepsilon_2). 


$$ 


 


Используя определение функции $\delta_m(t)$ через ряды, имеем очевидную 


оценку $0<\beta\le\linebreak 2(\varepsilon_1+\varepsilon_2)\tau^{2k}<1$.


Отсюда, прологарифмировав ее, получим (16) и 


$$ 


F_{\max}(r)=M^2(\lambda_2-\lambda_1)^2I^2_k(\sqrt{\lambda_2}r) 


I^{-2}_k(\sqrt{\lambda_2}R_2),\quad r\in[0,R_2]. 


$$ 


Оценка (15) следует из представления частного решения $\wt v(r)$. 


В общем случае по формулам Крамера решим систему уравнений (21),


причем


\begin{gather*}


\allowdisplaybreaks 


\Delta_k= 


\begin{vmatrix} 


a_{11} & a_{12} \\ a_{21} & a_{22} 


\end{vmatrix} 


=a_{22}a_{21}[\delta_{k_1}(\sigma_2)-\delta_{k_2}(\sigma_2)];\\


%


\Delta_{1,k}= 


\begin{vmatrix} 


2(\varepsilon_1+\varepsilon_2) & a_{12} \\ 


M^2(\lambda_2-\lambda_1)^2 & a_{22} 


\end{vmatrix} 


;\quad \Delta_{2,k}= 


\begin{vmatrix} 


a_{11} & 2(\varepsilon_1+\varepsilon_2) \\ 


a_{21} & M^2(\lambda_2-\lambda_1)^2 


\end{vmatrix} 


. 


\end{gather*} 


Введем обозначение $I^2_{k_i}(\sqrt{\lambda_2}r)=a_i(r)$, тогда 


\begin{multline*} 


F_{\max}(r)=\frac{\Delta_{1,k}}{\Delta_k} 


I^2_{k_1}(\sqrt{\lambda_2}r)+\frac{\Delta_{2,k}}{\Delta_k} 


I^2_{k_2}(\sqrt{\lambda_2}r) = \\ 


= \frac{\dfrac{a_1(r)}{a_{21}}-\dfrac{a_2(r)}{a_{22}}} 


{\delta^2_{k_1}(\sigma_2)-\delta^2_{k_2}(\sigma_2)} 


2(\varepsilon_1+\varepsilon_2)+ 


\frac{\dfrac{a_{11}a_2(r)}{a_{21}a_{22}}- 


\dfrac{a_{12}a_1(r)}{a_{22}a_{21}}} 


{\delta^2_{k_1}(\sigma_2)-\delta^2_{k_2}(\sigma_2)} 


M^2(\lambda_2-\lambda_1)^2. 


\end{multline*} 


Обе суммы монотонно возрастают по $k_1$ и монотонно убывают по $k_2$, 


что можно проверить с помощью дифференцирования или воспользоваться 


свойствами функции $\delta_k(t)$. Кроме того, 


$$ 


\lim_{k_1,k_2\to\left[\frac{1}{2}\ln2(\varepsilon_1+\varepsilon_2) 


-\ln\beta\right]/\ln\tau}\wt v\le M^2(\lambda_2-\lambda_1)^2 


\frac{I^2_k(\sqrt{\lambda_2}r)}{I^2_k(\sqrt{\lambda_2}R_2)}. 


$$ 


Отсюда и из представления общего решения через (5) следует оценка (15). 


Из неравенства (16) следует, что $k\to\infty$, если 


$\varepsilon_1+\varepsilon_2\to 0$, и тогда для больших 


$k$ 


имеем


$I_k(\sqrt{\lambda_2}r)/I_k(\sqrt{\lambda_2}R_2)\cong r^k/R^k_2$. 


Таким образом, для $r<R_2$ правая часть (15) стремится к нулю, 


и решение краевой задачи (1), (2) единственно в указанном выше классе. 


\end{pf} 


 


{\bf 3.} Используя дифференциальные свойства частного решения 


$v(r,\theta)$ 


по аргументам $r$ и $\theta$, можно получить более точные оценки, чем 


(18). Например, если вместо условия (15) потребовать выполнения условия 


\begin{equation} 


\frac{2}{\pi}\int^{2\pi}_0\bigg( 


\frac{\partial^m v}{\partial\theta^m}\bigg|_{r=R_2} 


\bigg)^2d\theta\le M^2, 


\end{equation} 


то для первого случая ($\lambda_1\ne\lambda_2$) получим оценку 


\begin{equation} 


\frac{2}{\pi}\int^{2\pi}_0u^2(r,\theta)d\theta\le2\bigg\{ 


\varepsilon_2+\frac{M^2(\lambda_2-\lambda_1)^2}{k^{2m}} 


\frac{I^2_k(\sqrt{\lambda_2}r)}{I^2_k(\sqrt{\lambda_2}R_2)}\bigg\}, 


\end{equation} 


где $k$ удовлетворяет неравенству 


\begin{equation} 


0<(1/2)\ln(\varepsilon_1+\varepsilon_2)+m\,\ln k+k\,\ln\tau-\ln\beta. 


\end{equation} 


 


Доказательство аналогично доказательству оценки (15), если вместо условия 


(14) использовать условие (22). Следовательно, имеет место 


 


\begin{cor} 


Пусть функция $\wt v(r,\theta)$\; $m$ раз дифференцируема по аргументу 


$\theta$ и выполняются 


неравенства 


$\|\varphi_i(\theta,R_i)\|_{L_2[0,2\pi]}\le\varepsilon_i$ ($i=1,2$) и 


(22). Тогда для разных корней уравнения (3) имеет 


место оценка (23), где $k$ определяется из (24), а для кратных 


корней --- оценка 


$$ 


\frac{2}{\pi}\int^{2\pi}_0u^2(r,\theta)d\theta\le2\bigg\{ 


\varepsilon_1+k^{-m}\bigg[ 


\frac{r I'_k(\sqrt{\lambda_2}r)}{R_2I'_k(\sqrt{\lambda_2}R_2)} 


\bigg]^2M^2\bigg\}, 


$$ 


где $k$ определяется из уравнения 


$$ 


M\frac{R_1}{R_2} 


\frac{I'_k(\sqrt{\lambda_2}R_1)}{I'_k(\sqrt{\lambda_2}R_2)} 


=k^m\sqrt{2(\varepsilon_1+\varepsilon_2)}\,. 


$$ 


\end{cor} 


 


Введем обозначения 


\begin{gather*} 


\frac{2}{\pi}\int^{2\pi}_0\biggl(\frac{\partial v}{\partial r} 


\bigg|_{r=R_2}\biggr)^2d\theta= 


\frac{1}{(\lambda-\lambda_2)^2}\sum^\infty_{k=0} 


A^2_k\Psi_k(r),\\ 


%


\Psi_k(r)=\frac{k^2}{r^2}I^2_k(\sqrt{\lambda_2}r)+2 


\frac{\sqrt{\lambda_2}}{r}k I_k(\sqrt{\lambda_2}r) 


I_{k+1}(\sqrt{\lambda_2}r)+\lambda_2I^2_{k+1} 


(\sqrt{\lambda_2}r). 


\end{gather*} 


По аналогии с предыдущими рассуждениями доказывается 


 


\begin{cor} 


Пусть выполняются условия 


$\|\varphi_i(\theta,R_i)\|_{L_2[0,2\pi]}\le\varepsilon_i$ ($i=1,2$), 


(4) и 


$$ 


\frac{2}{\pi}\int^{2\pi}_0\biggl(\frac{\partial v}{\partial r} 


\bigg|_{r=R_2}\biggr)^2d\theta\le M^2. 


$$ 


Тогда для $\lambda_1\ne\lambda_2$ имеет место оценка 


$$ 


\frac{2}{\pi}\int^{2\pi}_0 u^2(r,\theta)d\theta\le2\bigg\{ 


\varepsilon_1+ 


\frac{I^2_k(\sqrt{\lambda_2}r)}{\Psi_k(R_2)}M^2 


(\lambda_2-\lambda_1)^2\bigg\}, 


$$ 


где параметр $k$ определяется из уравнения 


$M^2\dfrac{I^2_k(\sqrt{\lambda_2}R_1)}{\Psi_k(R_2)}= 


2(\varepsilon_1+\varepsilon_2)$. 


\end{cor} 


 


\begin{note} 


Условия следствия 1 при $k=1$ и 2 можно заменить условием (11). 


\end{note} 


 


Аналогично рассматриваются вопросы условной устойчивости для уравнения 


(1) в пространствах с отличной от нуля кривизной \cite{3}. 


\medskip 


 


{\bf 4.} Рассмотрим случай двухконтурной задачи, когда внешняя граница 


$\Gamma_2$ --- окружность, а внутренняя задается уравнением 


$r=R_1(\theta)$, $r>0$. Обозначим 


определитель краевой задачи (2${}_4$), когда контуры $\Gamma_1$ и 


$\Gamma_2$ задаются соответственно уравнениями $r=R_1(\theta)$ и 


$r=R_2$, через


$$ 


\Delta_k(\theta)= 


\begin{vmatrix} 


I_m(\sqrt{\lambda_1}\,R_1(\theta)) &  


I_m(\sqrt{\lambda_2}\,R_1(\theta)) \\ 


I_m(\sqrt{\lambda_1}\,R_2) & I_m(\sqrt{\lambda_2}\,R_2) 


\end{vmatrix} 


. 


$$ 


 


Обозначим через $R_{10}$ радиус максимальной концентрической сферы, 


которую можно вписать в область, ограниченную кривой $r=R_1(\theta)$. 


 


\begin{lemnonum} 


Имеет место оценка 


$$ 


\Delta_k(\theta)>\Delta_{k0},


$$ 


где $\Delta_{k0}$ --- определитель двухконтурной задачи 


$$ 


u(r,\theta)\big|_{r=R_{10}}=\varphi_{10}(x),\quad 


u(r,\theta)\big|_{r=R_2}=\varphi_2(x). 


$$ 


\end{lemnonum} 


 


\begin{pf} 


Перепишем в форме (8) определитель 


$$ 


\Delta_k(\theta)=I_k(\sqrt{\lambda_1}\,R_2) 


I_k(\sqrt{\lambda_2}\,R_2)\bigg[ 


\frac{I_k(\sqrt{\lambda_1}\,R_1(\theta))} 


{I_k(\sqrt{\lambda_1}\,R_2)}- 


\frac{I_k(\sqrt{\lambda_2}\,R_1(\theta))} 


{I_k(\sqrt{\lambda_2}\,R_2)}\bigg] 


$$ 


и вычислим производную 


$$ 


\Delta'_k(\theta)=\Gamma_k\bigg[ 


\frac{I'_k(\sqrt{\lambda_1}\,R_1(\theta)) 


\sqrt{\lambda_1}\,R'_1(\theta)} 


{I_k(\sqrt{\lambda_1}R_2)}- 


\frac{I'_k(\sqrt{\lambda_2}\,R_1(\theta)) 


\sqrt{\lambda_2}\,R'_1(\theta)} 


{I_k(\sqrt{\lambda_2}\,R_2)}\bigg], 


$$ 


где $\Gamma_k=I_k(\sqrt{\lambda_1}\,R_2)/I_k(\sqrt{\lambda_2}\,R_2)$. 


Используя формулу дифференцирования функций Бесселя, 


для определителя $\Delta_k(\theta)$ получим дифференциальное уравнение 


первого порядка 


\begin{equation} 


\Delta'_k(\theta)=\Gamma_k 


\frac{k R'_1(\theta)}{R_1(\theta)}\Delta_k(\theta)+ 


g_k(\theta)R'_1(\theta), 


\end{equation} 


где 


$$ 


g_k(\theta)=\Gamma_k\bigg[ 


\frac{I_{k+1}(\sqrt{\lambda_1}\,R_1(\theta))} 


{I_k(\sqrt{\lambda_1}\,R_2)}\sqrt{\lambda_1}- 


\frac{I_{k+1}(\sqrt{\lambda_2}\,R_1(\theta))} 


{I_k(\sqrt{\lambda_2}\,R_2)}\sqrt{\lambda_2}\bigg], 


$$ 


общее решение которого запишем в виде 


$$ 


\Delta_k(\theta)=R^k_1(\theta)\Big[C_1+\int R^{-k}_1(\tau) 


g_k(\tau)R'_1(\tau)d\tau\Big], 


$$ 


$C_1$ --- произвольная константа. Нетрудно подсчитать, что 


$g_k(t)\ge0$, 


тогда вместо 


(25) получим неравенство $\Delta_k(\theta)>R^k_1(\theta)C_1$. Положив 


$C_1=R^{-k}_{10}\Delta_{k0}$ и учитывая, что 


$R_{10}=\min{}R_1(\theta)$, $\theta\in[0,2\pi]$, получим неравенство 


(20).


\end{pf} 





Обозначим $\tau_0=R_1/R_{10}$, тогда, следуя методике работ 


\cite{2}, \cite{3}, доказывается теорема существования и оценки условной 


устойчивости типа (15), (17) для двухконтурной задачи, если контуры 


$\Gamma_1$ и $\Gamma_2$ задаются соответственно уравнениями 


$r=R_1(\theta)$, $r=R_2$. 


 


\begin{thebibliography}{9} 


 


\bibitem{1} 


Бугир М.К. {\em Связь неколеблемости решений с некорректными задачами 


для уравнений эллиптического типа четвертого порядка} // Докл. РАН. -- 


1991. -- Т.\,319. -- \No\,5. -- С.\,1053--1056. 


\bibitem{2} 


Бугир М.К. {\em Некорректные краевые задачи для полигармонического 


уравнения $n$-го порядка} // Изв. вузов. Математика. -- 1994. -- 


\No\,10. -- С.\,19--25. 


\bibitem{3} 


Бугир М.К. {\em Краевые задачи для уравнений с частными производными на 


многообразиях постоянной кривизны} // Препринт. Ин-т матем. АН Украины. 


-- Киев, 1993. -- \No\,93. -- 58 с. 


\bibitem{4} 


Лаврентьев М.М., Романов В.Г., Шишатский С.П. {\em Некорректные задачи 


математической физики и анализа}. -- М.: Наука, 1980. -- 286 с. 


\bibitem{5} 


Векуа И.Н. {\em Новые методы решения эллиптических уравнений}. -- 


М.--Л.: ГИТТЛ, 1948. -- 296 с. 


\bibitem{6} 


Пташнык Б.И., Фиголь В.В., Штабалюк П.И. {\em Разрешимость, устойчивость 


и регуляризация многоточечной задачи для гиперболических уравнений} // 


Тр. Львов. матем. о-ва. -- 1991. -- Т.\,3. -- С.\,16--32. 


\bibitem{7} 


Курант Р. {\em Уравнения с частными производными}. -- М.: Мир, 1964. -- 


830 с. 


\bibitem{8} 


Камке Э. {\em Справочник по обыкновенным дифференциальным уравнениям}. 


-- М.: Наука, 1971. -- 576 с. 


 


\end{thebibliography} 


 


\vskip 2cm 


 


\postup{\it Тернопольская Академия}{$\qquad$\it Поступили} 


\postup{\it народного хозяйства}{{\it первый вариант} 03.03.1995} 


\postup{}{{\it окончательный вариант} 18.07.1996} 


\label{end} 


\end{document} 


